\documentclass[10pt]{article}
\usepackage[left=0.75in,right=0.75in,top=0.6in,bottom=0.6in]{geometry}
\usepackage{amsmath,booktabs,titlesec,parskip,microtype,enumitem,times,multicol}
\setlength{\parskip}{3pt}
\setlength{\parindent}{0pt}
\setlength{\columnsep}{0.25in}
\titleformat{\section}{\normalsize\bfseries}{}{0em}{}
\titleformat{\subsection}{\small\bfseries}{}{0em}{}
\titlespacing{\section}{0pt}{6pt}{2pt}
\titlespacing{\subsection}{0pt}{4pt}{1pt}
\title{\textbf{\large CRAFT: A Self-Building Agent Framework That Learns to Configure Telecom Networks}}
\author{Anonymous Authors}
\date{}
\begin{document}
\maketitle
\noindent\textbf{Abstract} - Configuring a telecom network site today is manual, repetitive, and operator-specific - every new operator restarts from zero with no learning, no reuse, and no carry-forward. \textbf{CRAFT} (Configuration RAN Auto-generation Framework for Telecom) solves this with six AI agents that \emph{learn} rules from past deployments, \emph{reason} with 3GPP standards, \emph{discover} output templates automatically, \emph{stress-test} every rule, \emph{build} missing tools on the fly, and \emph{improve} from every correction. Tested on real Vodafone and Telus deployments, CRAFT achieves \textbf{95.2\% field-level accuracy}, reduces generation time from 2-3 days to 4 minutes, and reaches \textbf{78.6\% rule coverage on a brand-new operator with zero prior exposure}.
\vspace{4pt}\hrule\vspace{4pt}
\begin{multicols}{2}
\section{Introduction}
\subsection*{The Problem: Every Operator is a New Project}
Picture an engineer on a Monday morning. She opens Vodafone's CIQ - 1,190 rows across 6 sheets - then a 20-sheet LLD, then a rule book with 81 handcrafted rules. She writes Java code. By Friday the site is done. She does this 50 times that month.
Then Telus arrives. Different files, different rules, different outputs. The entire codebase must be rewritten. Every operator is a new software project with no carry-over.
\subsection*{Why Existing Approaches Fall Short}
Recent LLM-based approaches~\cite{lira2024largelanguagemodelszero,provvedi2026intentllm,wang2026llmpowered} reduce some manual effort but still require per-operator engineering. They cannot \emph{learn} from deployed configurations, \emph{reason} with telecom standards, or \emph{build} missing tools. Every new operator remains a blank slate.
CRAFT's insight: the knowledge to configure any operator already exists - embedded in past deployments, encoded in 3GPP standards, and latent in cross-operator patterns. Six specialized agents extract, organize, and apply this knowledge automatically.
\section{The CRAFT Solution}
CRAFT operates through five sequential stages, with a sixth agent running continuously to feed improvements back.
\subsection*{Stage 1: Learning from History (ISA)}
The \textbf{Inverse Synthesis Agent (ISA)} mines past deployed configurations and reverse-engineers the rules that produced them. Identical values across sites become default rules; varying values are traced to CIQ or LLD sources. On Vodafone, ISA extracts \textbf{45 of 81 rules} with confidence above 0.8. The remaining 36 are flagged for DGRA - the system knows the limits of its own knowledge.
\subsection*{Stage 2: Reasoning with Domain Knowledge (DGRA)}
The \textbf{Domain-Grounded Rule Agent (DGRA)} is a fine-tuned LLM with retrieval access to 3GPP standards (TS 38.331, TS 36.331), Samsung specifications, and a cross-operator rule library. It resolves three classes ISA cannot:
\begin{itemize}[leftmargin=1em,itemsep=1pt,topsep=1pt]
  \item \textbf{Multi-source derivation} - values requiring chained CIQ and LLD lookups (e.g., \texttt{site\_config} determines which LLD sheet to read).
  \item \textbf{Cross-operator transfer} - vector similarity finds analogous rules from other operators. On Telus (never seen before), DGRA achieves \textbf{78.6\% rule coverage} purely from Vodafone patterns.
  \item \textbf{3GPP grounding} - radio parameter rules validated against telecom standards before acceptance.
\end{itemize}
Together, ISA and DGRA cover \textbf{91.4\% of Vodafone rules} before any template is touched.
\subsection*{Stage 3: Selecting Output Templates (TDA)}
The \textbf{Template Discovery Agent} determines output structure from input data: \texttt{site\_config} identifies the NE type, \texttt{technology} selects cell sections, vDU count sets template count. Previously hand-coded per operator; TDA discovers it automatically.
\subsection*{Stage 4: Testing and Tool Building (MTA + PVS)}
The \textbf{Mutation Testing Agent} stress-tests every rule by injecting boundary values, null fields, and conflicting inputs. Failed rules are hardened with validation guards. CRAFT is the first telecom configuration system that actively \emph{tests its own rules}.
The \textbf{Plan-Verify-Synthesize (PVS) loop} builds missing tools: generates specifications, writes code via LLM, tests in sandbox, registers validated tools. \textbf{5 of 6} first-attempt syntheses succeeded. A confidence gate governs execution: $>$0.9 auto-fills; 0.7-0.9 logged; 0.5-0.7 human review; $<$0.5 skipped.
\subsection*{Stage 5: Learning from Corrections (FLA)}
The \textbf{Feedback Learning Agent} captures corrections, identifies root causes, updates rules, and propagates improvements across all operators - a fix for Vodafone immediately improves Telus.
\section{Results}
Evaluated on two operators with entirely different formats and no shared templates.
\textbf{Vodafone:} CIQ (1,190$\times$88, 6 sheets) + LLD (20 sheets) + 81 rules. \textbf{Telus:} POD Placement + Grow Rules + BPI + MME/RU.
\begin{center}
\renewcommand{\arraystretch}{1.2}
\small
\begin{tabular}{lcc}
\toprule
\textbf{Metric} & \textbf{Vodafone} & \textbf{Telus (zero-shot)} \\
\midrule
Rules by ISA & 45/81 & - \\
Rules by DGRA & 28/36 & 78.6\% \\
Total coverage & 91.4\% & 78.6\% \\
Field accuracy & 95.2\% & in prog. \\
Tool synthesis & 5/6 & - \\
Gen. time & 4 min & 4 min \\
Manual baseline & 2-3 days & 2-3 days \\
\bottomrule
\end{tabular}
\end{center}
The Telus result is most significant. Without prior exposure, CRAFT achieved 78.6\% rule coverage using cross-operator learning and 3GPP grounding - the compounding effect in action.
\section{Impact}
\textbf{Faster onboarding.} 6-12 months to 2-4 weeks (85\% reduction). \textbf{Error reduction.} 5-8\% to below 0.5\%. \textbf{Compounding knowledge.} Each deployment enriches the cross-operator rule base. \textbf{AI-native RAN.} Positions Samsung at the forefront of autonomous network management.
\section{Conclusion}
Telecom configuration is a knowledge problem, not a coding problem. CRAFT extracts, validates, applies, and improves this knowledge through six agents in a self-reinforcing pipeline. Every new operator makes all previous ones more accurate - a compounding advantage that grows with scale.
\end{multicols}
\begin{thebibliography}{3}
\small
\bibitem{lira2024largelanguagemodelszero} O.~G. Lira et al., ``Large language models for zero touch network configuration management,'' arXiv:2408.13298, 2024.
\bibitem{provvedi2026intentllm} C.~Provvedi et al., ``Intent-LLM: A framework for automated network configuration through code generation,'' IEEE TCCN, 2026.
\bibitem{wang2026llmpowered} J.~Wang et al., ``LLM-powered intent-driven configuration generation for multi-vendor networks,'' IEEE TNSM, vol.~23, pp.~3537-3555, 2026.
\end{thebibliography}
\end{document}
